\section*{अध्याय \ref{ch:g8:powers} --- घातें}

\begin{solution}{exo:g8:powers:1}
$2^6 = 64$; \quad $3^3 = 27$; \quad $10^5 = 100\,000$; \quad
$1^{100} = 1$; \quad $0.1^2 = 0.01$; \quad $6^0 = 1$.
\end{solution}

\begin{solution}{exo:g8:powers:2}
$(-3)^2 = 9$; \quad $-3^2 = -9$; \quad $(-1)^{15} = -1$; \quad
$(-5)^3 = -125$; \quad $\left(\frac23\right)^2 = \frac49$.
\end{solution}

\begin{solution}{exo:g8:powers:3}
$7^4 \times 7^5 = 7^9$; \quad
$\dfrac{2^9}{2^3} = 2^6$; \quad
$\left(10^3\right)^4 = 10^{12}$; \quad
$5^6 \times 5^{-2} = 5^4$; \quad
$3^4 \times 7^4 = 21^4$ (एक ही \omterm{def:g8:powers:def}{घातांक}: आधार गुणा कर दो)।
\end{solution}

\begin{solution}{exo:g8:powers:4}
$10^{-2} = 0.01$; \quad
$4 \times 10^3 = 4000$; \quad
$2.5 \times 10^{-4} = 0.00025$; \quad
$10^0 = 1$.
\end{solution}

\begin{solution}{exo:g8:powers:5}
$10^5$; \quad $10^{-1}$; \quad $10^{10}$; \quad $10^{-6}$.
\end{solution}

\begin{solution}{exo:g8:powers:6}
$\dfrac{10^7 \times 10^{-3}}{10^2} = \dfrac{10^4}{10^2} = 10^2 =
100$.

$\dfrac{2^5 \times 2^4}{2^6} = \dfrac{2^9}{2^6} = 2^3 = 8$.

$\left(2^2\right)^3 \times 2^{-4} = 2^6 \times 2^{-4} = 2^2 = 4$.
\end{solution}

\begin{solution}{exo:g8:powers:7}
$2^3 \times 2^4 = 2^{12}$: \emph{ग़लत} --- \omterm{def:g8:powers:def}{घातांक} जुड़ते हैं:
$2^7$।

$5^2 + 5^3 = 5^5$: \emph{ग़लत} --- योग के लिए कोई नियम नहीं है:
$25 + 125 = 150$, जबकि $5^5 = 3125$।

$(3^2)^4 = 3^8$: \emph{सही}।

$10^3 \times 10^3 = 100^3$: \emph{सही} --- दोनों $10^6$ के बराबर
हैं (बाएँ: $10^{3+3}$; दाएँ: $(10^2)^3$)।
\end{solution}

\begin{solution}{exo:g8:powers:8}
चौथे दिन के नए लोग: तीसरे दिन ख़बर पाने वाले $3^3 = 27$ लोगों में
से हर एक तीन और को बताता है: $3^4 = 81$। चौथे दिन के अंत तक जानने
वाले: $3 + 9 + 27 + 81 = 120$ लोग।
\end{solution}

\begin{solution}{exo:g8:powers:9}
\emph{1.} मोटाई: $2^{10} \times 10^{-4}$ m
$= 1024 \times 10^{-4}$ m $\approx 0.1$ m $= 10$ cm।

\emph{2.} $2^{42} \times 10^{-4} \approx 4.4 \times 10^{12} \times
10^{-4} = 4.4 \times 10^8$ m, जो $3.8 \times 10^8$ m से बड़ा है:
$42$ बार (काग़ज़ पर!) मोड़ने के बाद काग़ज़ का गट्ठर चंद्रमा से आगे
निकल जाता।
\end{solution}

\begin{solution}{exo:g8:powers:10}
$2^{10} = 1024$; $10^3 = 1000$; $3^6 = 729$; $5^4 = 625$। क्रम:
\[
5^4 < 3^6 < 10^3 < 2^{10} .
\]
\end{solution}

\begin{solution}{exo:g8:powers:11}
$2^{100} = \left(2^{10}\right)^{10} = 1024^{10}$ और
$10^{30} = \left(10^3\right)^{10} = 1000^{10}$। चूँकि
$1024 > 1000$ है, दोनों की दस-दस नक़लें गुणा करने पर असमिका बनी
रहती है: $2^{100} > 10^{30}$।
\end{solution}

\begin{solution}{pb:g8:powers:1}
\textbf{1.} दाने $1 = 2^0$ से शुरू होकर ख़ाने-दर-ख़ाने दुगुने होते
हैं: ख़ाने $k$ पर $2^{k-1}$ दाने हैं
(\cref{def:g8:powers:def})। ख़ाने $64$ पर $2^{63}$ हैं।

\textbf{2.} $S_1 = 1$, $S_2 = 1 + 2 = 3$, $S_3 = 3 + 4 = 7$,
$S_4 = 7 + 8 = 15$, $S_5 = 15 + 16 = 31$: हर बार $2$ की अगली घात
($2$, $4$, $8$, $16$, $32$) से एक कम। अटकल:
$S_n = 2^n - 1$।

\textbf{3.} $S_n$ के हर पद को दुगुना करने से हर घात एक क़दम ऊपर
सरक जाती है ($2 \times 2^{k} = 2^{k+1}$):
\begin{align*}
S_n &= 1 + 2 + 4 + \dots + 2^{n-1}, \\
2 \times S_n &= \phantom{1 + {}} 2 + 4 + \dots + 2^{n-1} + 2^n .
\end{align*}
पहली पंक्ति को दूसरी में से घटाने पर $2$ से $2^{n-1}$ तक का हर पद
दोनों में आता है और कट जाता है:
\[
2 \times S_n - S_n = 2^n - 1,
\qquad\text{यानी}\qquad
S_n = 2^n - 1 .
\]

\textbf{4.} कुल $S_{64} = 2^{64} - 1$ दाने हैं। एक पैंसठवें ख़ाने
पर $2^{64}$ दाने होते: पूरी बिसात पर ठीक उससे एक दाना कम है।

\textbf{5.} पहले आधे हिस्से पर $S_{32} = 2^{32} - 1$ दाने हैं।
पूरी बिसात पर $2^{64} - 1$ हैं, इसलिए दूसरे आधे हिस्से पर
\[
\left(2^{64} - 1\right) - \left(2^{32} - 1\right)
= 2^{64} - 2^{32}
= 2^{32} \times \left(2^{32} - 1\right)
\]
($2^{32}$ को बाहर निकालकर, और $2^{32} \times 2^{32} = 2^{64}$ का
इस्तेमाल करते हुए, \cref{thm:g8:powers:rules}): यानी ठीक पहले आधे
हिस्से का $2^{32}$ गुना।

\textbf{6.} \omterm{def:g8:powers:def}{घातांक} के नियमों से $2^{64} = 2^{4 + 60} =
2^4 \times \left(2^{10}\right)^6$। और $2^{10} = 1024 > 10^3$ है,
इसलिए
\[
2^{64} > 16 \times \left(10^3\right)^6 = 16 \times 10^{18}
= 1.6 \times 10^{19} .
\]

\textbf{7.} $1.6 \times 10^{19}$ से ज़्यादा दाने, हर एक
$5 \times 10^{-2}$ g का:
\[
1.6 \times 10^{19} \times 5 \times 10^{-2}
= 8 \times 10^{17} \text{ g} .
\]
$10^6$ g प्रति टन से भाग देने पर: $8 \times 10^{11}$ टन से ज़्यादा
--- यानी आठ सौ बिलियन टन।

\textbf{8.} $\dfrac{8 \times 10^{11}}{8 \times 10^{8}} = 10^3$:
राजा ने आज की पूरी दुनिया की कम से कम \emph{एक हज़ार साल} की फ़सल
का वादा कर डाला। (किंवदंती कहती है कि उसके सलाहकारों ने उसे यही
बताया था।)

\textbf{9.} $2^{19} = 2^9 \times 2^{10} = 512 \times 1024 =
524\,288 < 10^6$, जबकि $2^{20} = \left(2^{10}\right)^2 =
1024^2 = 1\,048\,576 > 10^6$। इसलिए सबसे छोटा \omterm{def:g8:powers:def}{घातांक} $n = 20$ है:
बीस बार दुगुना करने पर दस लाख पार हो जाता है।

\textbf{10.} दिन $n$ की तनख़्वाह $2^{n-1}$ सेंट है (पहला दिन:
$2^0 = 1$)। अब
\[
2^{26} = 2^6 \times \left(2^{10}\right)^2
= 64 \times 1\,048\,576 = 67\,108\,864 < 10^8,
\]
\[
2^{27} = 2 \times 2^{26} = 134\,217\,728 > 10^8 .
\]
इसलिए उस दिन की तनख़्वाह पहली बार $10^8$ सेंट से ऊपर तब निकलती है
जब $n - 1 = 27$ हो: यानी अट्ठाईसवें दिन।

\textbf{11.} $21 = 16 + 4 + 1$: बाट $16$, $4$ और $1$ g के।
$27 = 16 + 8 + 2 + 1$: बाट $16$, $8$, $2$ और $1$ g के।

\textbf{12.} सवाल 3 से बाट $1, 2, 4, 8$ मिलकर
$S_4 = 2^4 - 1 = 15$ g बनाते हैं। इसलिए $16$ g वाले बाट के बिना
व्यापारी $15$ g से आगे जा ही नहीं सकती: $16$ g या उससे ज़्यादा के
सामान में वह बाट लगना ही है। और $15$ g या उससे कम के सामान में वह
लग नहीं सकता, क्योंकि अकेला $16$ g वाला बाट ही सामान से भारी है।
इसलिए सबसे बड़े बाट का चुनाव \emph{तय} है। अब बचता है ज़्यादा से
ज़्यादा $15$ g का सामान, जिसे $1, 2, 4, 8$ से बनाना है --- और वही
तर्क दोहरा जाता है: $8$ तय है (बचे सामान के $\geq 8$ होने पर लगेगा,
वरना नहीं, क्योंकि $1 + 2 + 4 = 7$), फिर $4$ (क्योंकि
$1 + 2 = 3$), फिर $2$, फिर $1$।

\textbf{13.} इन तय चुनावों पर चलते हुए, हर चरण के बाद बचा हुआ
सामान बचे हुए बाटों के कुल से ज़्यादा नहीं होता, इसलिए यह सिलसिला
\omterm{def:g3:division:remainder}{शेषफल} $0$ पर ख़त्म होता है: $1$ से $31$ तक का हर सामान तोल लिया
जाता है। और चूँकि रास्ते का हर चुनाव तय था, इसलिए बाटों का कोई और
चुनाव उसी सामान तक नहीं पहुँच सकता: $1$ से $31$ तक की हर संख्या का
अलग-अलग $2$ की घातों के योग वाला लेखन होता है और वह \emph{अकेला}
होता है।

\textbf{14.} छहों बाट मिलकर $S_6 = 2^6 - 1 = 63$ g बनाते हैं, और
वही तय-चुनाव वाला तर्क $1$ से $63$ g तक के हर सामान को ढक लेता है।
$45$ के लिए: सामान $\geq 32$ है, इसलिए $32$ लगाओ; बचा $13 < 16$,
इसलिए $16$ छोड़ो; $8$ लगाओ (बचा $5$), $4$ लगाओ (बचा $1$), $2$
छोड़ो, $1$ लगाओ:
\[
45 = 32 + 8 + 4 + 1 = 2^5 + 2^3 + 2^2 + 2^0 .
\]

\textbf{15.} $64$ ख़ानों में से हर एक पर $0$ या $1$ चुनना असल में
यह चुनना है कि योग में $2^0, 2^1, \dots, 2^{63}$ में से कौन-कौन सी
घातें डालनी हैं --- ठीक वैसे ही जैसे व्यापारी $64$ बाटों से तोलती।
सबसे छोटी रखी जा सकने वाली संख्या $0$ है (सारे ख़ाने $0$ पर); और
सबसे बड़ी \emph{सारी} घातों का योग है, जो राजा वाला कुल ही है:
$S_{64} = 2^{64} - 1$ (भाग I)। तय-चुनाव वाले तर्क से बीच की हर
पूर्ण संख्या ठीक एक बार बनती है: इसलिए $64$-बिट मशीन ठीक $0$ से
$2^{64} - 1$ तक की पूर्ण संख्याएँ ही रखती है।

\end{solution}
